\documentclass[11pt]{article}
\usepackage[margin=1in]{geometry}
\usepackage{graphicx}
\usepackage{amsmath}
\usepackage{listings}
\usepackage{xcolor}
\usepackage{hyperref}
\usepackage{float}

% Code styling
\lstset{
    basicstyle=\ttfamily\small,
    breaklines=true,
    frame=single,
    backgroundcolor=\color{gray!10},
    numbers=left,
    numberstyle=\tiny,
    language=Python
}

\title{Kalshi Implied Distribution: CPI Year-Over-Year Analysis}
\author{}
\date{}

\begin{document}

\maketitle

\section{Overview}

This project extracts and visualizes the implied probability distribution from Kalshi's KXCPIYOY (CPI Year-Over-Year) ladder market using a three-stage pipeline: market data extraction, monotonic enforcement, and distribution smoothing.

\section{Market Selection}

I chose the CPI Year-Over-Year ladder market because it offers dense strike coverage (typically 10-15 contracts) with substantial liquidity, making it ideal for distribution extraction. The market structure provides binary contracts at incrementing thresholds (e.g., ``Will CPI YoY be greater than or equal to 2.0\%?''), where each contract's price represents the market's cumulative probability $P(\text{CPI} \geq \text{strike})$.

\section{Methodology}

\subsection{Stage 1: Market Data Extraction}

I fetch orderbook data for all active contracts in the nearest expiration, computing a weighted micro-price for each strike. Instead of simple mid-prices, I use order book imbalance weighting:

\begin{equation*}
\text{Micro-Price} = \frac{\text{Bid} \times \text{Ask Size} + \text{Ask} \times \text{Bid Size}}{\text{Total Size}}
\end{equation*}

This captures true market sentiment by weighting toward the heavier side of the book, reducing noise from stale quotes.

\begin{figure}[H]
\centering
\includegraphics[width=0.9\textwidth]{step1_market_data.png}
\caption{Step 1: Enhanced Market Data weighted by Open Interest. The uncertainty cloud shows bid-ask spreads, while bubble size indicates open interest concentration.}
\label{fig:step1}
\end{figure}

\noindent \textbf{Code snippet - Micro-price calculation:}
\begin{lstlisting}
# Weighted micro-price using order book imbalance
total_size = yes_bid_size + yes_ask_size
if total_size > 0:
    micro_price = (yes_bid_price * yes_ask_size + 
                   yes_ask_price * yes_bid_size) / total_size
else:
    micro_price = (yes_bid_price + yes_ask_price) / 2
\end{lstlisting}

\subsection{Stage 2: Monotonic Enforcement}

Raw market prices often violate monotonicity due to liquidity fragmentation or arbitrage delays. I apply weighted isotonic regression with liquidity as sample weights:

\begin{equation*}
\min_{f} \sum_{i=1}^{n} w_i (y_i - f(x_i))^2 \quad \text{subject to } f(x_1) \geq f(x_2) \geq \cdots \geq f(x_n)
\end{equation*}
where $w_i = \text{Liquidity}_i + \epsilon$ ensures $P(\text{CPI} \geq X)$ decreases monotonically as $X$ increases. The weighting prioritizes high-liquidity strikes, allowing the algorithm to correct illiquid outliers while preserving signals from deep markets.

% INSERT FIGURE: Step 2 visualization (Isotonic Regression)
\begin{figure}[H]
\centering
\includegraphics[width=0.85\textwidth]{step2_isotonic.png}
\caption{Step 2: Weighted Isotonic Regression. The red curve enforces monotonicity while respecting liquidity-weighted market prices (gray bubbles).}
\label{fig:step2}
\end{figure}
\noindent \textbf{Code snippet - Weighted isotonic regression:}
\begin{lstlisting}
X = df['Strike'].values
y = df['Mid'].values
weights = df['Liquidity'].values + 1e-6

# Weighted isotonic regression (decreasing)
iso_reg = IsotonicRegression(increasing=False, out_of_bounds='clip')
y_monotonic = iso_reg.fit_transform(X, y, sample_weight=weights)
\end{lstlisting}

\subsection{Stage 3: Distribution Smoothing}

I convert the step function into a continuous probability density using cubic spline interpolation. Two critical refinements ensure validity:

First, I add anchor points at the tails ($P=1.0$ at $0\%$, $P=0.0$ at $\text{max\_strike} + 0.2\%$) to prevent extrapolation artifacts. Without these, the spline often hallucinates long right tails that inflate expected values.

Second, I normalize the extracted PDF to integrate to 1.0. The derivative of the CDF spline can produce negative noise or violate total probability; clipping negatives and normalizing ensures a valid distribution:

\begin{align}
\text{PDF}_{\text{raw}}(x) &= \frac{d}{dx}\text{CDF}_{\text{spline}}(x) \\
\text{PDF}_{\text{clipped}}(x) &= \max(0, \text{PDF}_{\text{raw}}(x)) \\
\text{PDF}_{\text{normalized}}(x) &= \frac{\text{PDF}_{\text{clipped}}(x)}{\int \text{PDF}_{\text{clipped}}(x) \, dx}
\end{align}

% INSERT FIGURE: Step 3 visualization (Final PDF)
\begin{figure}[H]
\centering
\includegraphics[width=0.9\textwidth]{step3_final_pdf.png}
\caption{Final Distribution: Implied Probability Density with normalized mass. The vertical dashed line indicates market expectation, and rug plot marks liquid strikes.}
\label{fig:step3}
\end{figure}

\noindent \textbf{Code snippet - PDF extraction and normalization:}
\begin{lstlisting}
# Extract and normalize PDF
x_grid = np.linspace(df['Strike'].min(), max_strike + 0.2, 500)
pdf_series = cdf_spline.derivative()(x_grid)
pdf_series = np.maximum(pdf_series, 0)
total_area = np.trapezoid(pdf_series, x_grid)
pdf_series = pdf_series / total_area

# Calculate statistics
expected_val = np.trapezoid(x_grid * pdf_series, x_grid)
\end{lstlisting}

\section{Results}

For the December 2025 expiration, the model extracted a tight distribution with an expected CPI YoY of approximately 2.6--2.8\%, with most probability mass concentrated within a 0.5\% range. The normalized distribution shows total probability mass of 1.00, confirming calibration.

\section{Technical Validation}

Calibration is verified through three checks: (1) monotonicity of cumulative probabilities, (2) total probability mass equals 1.0, and (3) expected value lies within the range of liquid strikes. The visualization includes a rug plot highlighting liquid contracts, showing that the model concentrates probability where market information is most reliable.

\section{Key Innovations}

The combination of micro-price weighting, liquidity-weighted isotonic regression, and tail-controlled spline fitting addresses the core challenge of ladder markets: converting noisy, discrete cumulative probabilities into a coherent continuous distribution. Unlike simple interpolation, this approach respects market microstructure while enforcing mathematical validity.

The code is modular and generalizable to any Kalshi ladder market, requiring only the series ticker change. All computations use standard scientific libraries (scipy, scikit-learn) for reproducibility.

\end{document}
